\begin{center}\textbf{Résumé}\end{center}
\begingroup
\fontsize{9.5}{10.7}\selectfont
\setlength{\parskip}{1.5pt plus .5pt minus .5pt}
\noindent
On construit un prolongement de la structure discrète du continu qui
relie ses coordonnées numériques \(\pi,\varphi,e,\alpha\) à l’incidence exceptionnelle, aux
algèbres d’opérateurs de vertex et à la dualité d’un secteur compact.
L’holographie modulaire triadique (HMT) part de \emph{All Possible Paths}
(APP) : le support \((\mathbb Z/9\mathbb Z)^2\), son graphe de Cayley additif
avec générateurs cardinaux \(\pm(1,0),\pm(0,1)\) et les évaluations somme
et produit, avec résidus et quotients conservés. La signature
\(\mathrm{TRIT}=\{-1,0,+1\}\) distingue l’orientation et le régime ; le
\emph{Topological Phase Kernel} (TPK) compose la sélection d’évaluation,
le transport et la mise à jour de mémoire pour prolonger des états compatibles.
Le même arbre d’histoires soutient la mesure cylindrique, le transport
des feuillets, le bilan orienté, l’incidence et les complexes de mémoire
avec leurs applications de raffinement.

Le registre dodécaphasique engendré \(K\) relie deux réalisations.
Son incidence détermine un drapeau de Witt qui, avec le recollement, la métrique
et la chambre déclarés, conduit à un voisin de Leech. Les théorèmes de
construction d’algèbres de vertex applicables à ce réseau identifient
les extensions d’ordres deux et trois à \(V^\natural\), dont le groupe
d’automorphismes est le Monstre ; les produits, la graduation et les
traces sont transportés au moyen d’isomorphismes choisis. Dans la représentation à douze
positions de \(A_5\), la projection de \(K\) sur le secteur
tridimensionnel est non nulle et détermine la réduction orthogonale
\(12\to11\to10\). L’orbite cyclique du registre permet d’identifier
des opérateurs avec des bornes de stabilité ; le bilan bilatéral fournit
la récupération de l’état et de la mémoire sous ses hypothèses isométriques.

% BEGIN S27_FRAMING
La comparaison des étoiles de \(B_K\), \(B_\varphi\) et \(B_{\rm APP}\) distingue
les signatures polarisées \((3,1,0)\), \((1,0,3)\) et \((1,1,2)\).
Le repère ordonné des quatre supports
\(H^-_\alpha,H_e,H_\varphi,H_{\rm APP}\) a un stabilisateur trivial
dans l’action de Witt sur les douze positions. La lecture numérique
conserve, en plus des résidus régionaux, les quotients entiers de
somme et de différence : leur composition satisfait une identité de cocycle,
et leurs réductions apportent des composantes algébriques supplémentaires à
l’évaluation affine des entrées. Ces résultats précisent les
fonctions complémentaires d’incidence, d’orientation et de mémoire de
retenue dans le prolongement commun.

% END S27_FRAMING
La réalisation de Weyl de la phase engendre les algèbres matricielles de chaque
profondeur ; les inclusions de fibre complète produisent une algèbre
uniformément hyperfinie et un facteur tracial de type \(II_1\).
La section d’action et la paire angulaire fixent une ellipse orientée et un
rayon normalisé unique. L’échange des coordonnées entières et
l’inversion du rayon conservent sa forme quadratique et entrelacent
unitairement les Hamiltoniens autoadjoints, y compris leurs domaines.
Le transport de la relation résiduelle garantit la covariance du changement
de section. Un morphisme d’écrans coordonne cette dualité avec la réduction
dirigée par \(K\) et le quotient isotrope réticulaire \(26\to24\).

La composition des trajectoires à bord conservé définit un cocycle
additif d’action. Son prolongement vers la théorie M spécifie le
codomaine des champs à onze dimensions et les applications de réalisation qui
doivent conserver l’action, la graduation et la dualité. Sous les entrelacements
et domaines déclarés, l’identité de supercharge est transportée à l’image
isométrique. Les conditions d’identification physique complète sont
distinguées de ces résultats algébriques et opératoriels.
\par\medskip

% BEGIN R27_FRAMING_SUMMARY_ES
L’arithmétique qui précède ces réalisations inclut les lois du
défaut APP sous changement de module et de dimension, le sous-réseau invariant
d’indice deux et le codage nonadique avec retenue. L’émetteur détermine
un alphabet de \(42\) blocs et des raccordements régionaux explicites ; le registre
ordonné conserve des corrélations qu’il distingue de son histogramme. La
corrélation quadratique démontre l’identité du système de coordonnées de bord,
et la sélection sépare la saturation, la neutralité duale et l’orientation. L’identité
rationnelle de niveau trois établit en outre une divisibilité
uniforme optimale par \(3^9\) dans l’anneau des séries formelles.
% END R27_FRAMING_SUMMARY_ES

\noindent\textbf{Mots-clés :} holographie modulaire triadique ; structure
discrète du continu ; registre dodécaphasique ; algèbres d’opérateurs de
vertex ; Moonshine ; Monstre ; dualité T ; représentation de Weyl ;
écrans dimensionnels ; théorie M.
\par\endgroup
